\documentclass[10pt]{article}
\usepackage[margin=1in]{geometry}
\usepackage[T1]{fontenc}
\usepackage{lmodern,microtype,amsmath,amssymb,amsthm}
\theoremstyle{definition}
\newtheorem{assumption}{Assumption}
\begin{document}
\section*{Nonconvex Optimization with Mixed Relative Noise}
We consider a differentiable, possibly nonconvex objective $f:E\to\mathbb{R}$ on a separable real Hilbert space $E$. Its gradient is Lipschitz continuous, and its value at a given initial point $x^0\in E$ exceeds a global lower bound by a bounded amount. Each oracle call returns an unbiased stochastic gradient at the queried point, with a relative-plus-additive variance bound and a mixed relative-and-distance bound on noise differences. Consecutive pairs of calls share a sample; different pairs use independent samples. For a prescribed $\varepsilon>0$, the goal is a measurable output $\bar{x}\in E$ with integrable gradient norm satisfying $\mathbb{E}\|\nabla f(\bar{x})\|\leq\varepsilon$, where expectation includes all algorithm and oracle randomness.

The underlying optimization problem is
\[
\inf_{x\in E} f(x),
\]
where the norm on $E$ is induced by its inner product. Equip $E$ with its Borel sigma-algebra, and let $\Omega$ be a measurable sample space. The public inputs are $E$, $\Omega$, $x^0$, and the parameters
\[
L>0,\qquad M,r,\sigma,\Delta\geq 0,\qquad \varepsilon>0.
\]
The objective and its stochastic oracle satisfy the following conditions.

\begin{assumption}[Smoothness and initial gap]
The function $f$ is differentiable everywhere, and for every $x,y\in E$,
\[
\begin{aligned}
\|\nabla f(x)-\nabla f(y)\| &\leq L\|x-y\|,\\
f(x^0)-f(x) &\leq \Delta.
\end{aligned}
\]
Thus $f(x^0)-\Delta$ is a global lower bound, which need not be attained.
\end{assumption}

\begin{assumption}[Unbiasedness and relative variance]
There exist a probability measure $p$ on $\Omega$ and a jointly measurable map $g:E\times\Omega\to E$. Define the centered noise by
\[
 e(x,\omega)=g(x,\omega)-\nabla f(x).
\]
For every $x\in E$, the function $\omega\mapsto\|e(x,\omega)\|^2$ is integrable under $p$, and
\[
\begin{aligned}
\int_\Omega g(x,\omega)\,dp(\omega) &= \nabla f(x),\\
\int_\Omega \|e(x,\omega)\|^2\,dp(\omega)
&\leq r\|\nabla f(x)\|^2+\sigma^2.
\end{aligned}
\]
The vector integral is a Bochner integral; its integrability follows from joint measurability and the integrable second moment of the centered noise.
\end{assumption}

\begin{assumption}[Mixed noise regularity]
For every $x,y\in E$, noise differences evaluated using the same sample satisfy
\[
\begin{aligned}
&\int_\Omega \|e(x,\omega)-e(y,\omega)\|^2\,dp(\omega)\\
&\qquad\leq r\bigl(\|\nabla f(x)\|^2+\|\nabla f(y)\|^2\bigr)
+M^2\|x-y\|^2.
\end{aligned}
\]
\end{assumption}

\begin{assumption}[Paired-sample oracle access]
Let $P$ be a probability measure on $\Omega^{\mathbb{N}}$ under which the coordinate variables $(\xi_j)_{j\geq 0}$ are measurable, mutually independent, and each has law $p$. At zero-based call index $k$, the single oracle, with ID $0$, returns
\[
 g\bigl(x,\xi_{\lfloor k/2\rfloor}\bigr)\in E
\]
at the query point $x\in E$. Thus calls $2j$ and $2j+1$ share $\xi_j$. Each evaluation costs one oracle call, including the second evaluation in a pair. Internal algorithm randomness is independent of the entire sample sequence.
\end{assumption}

Queries may be chosen adaptively anywhere in $E$, including choosing the second point in a pair after observing the first reply. Unbiasedness above is a fixed-point statement under $p$; it does not assert conditional unbiasedness for this second adaptive query. The algorithm accesses $f$ through the oracle and is selected using only the public inputs; $f$, $g$, $p$, and $P$ are not supplied to it. A call budget $K\in\mathbb{N}$ may depend on the full public parameter tuple.

Within this budget, the algorithm must return a measurable point $\bar{x}\in E$. Writing $\mathbb{P}$ for the joint law of its independent internal randomness and the oracle samples, the exact success requirement is
\[
\|\nabla f(\bar{x})\|\in L^1(\mathbb{P}),
\qquad
\int \|\nabla f(\bar{x})\|\,d\mathbb{P}\leq\varepsilon.
\]
\end{document}
